\chapter{Related Work on \gls{rl} for Compiler Optimization}
\label{ch:related-work}

\section{Introduction}

The application of reinforcement learning to compiler optimization sits at the intersection of two rich research traditions: decades of work on compiler transformation frameworks, and the more recent emergence of deep \gls{rl} as a practical tool for sequential decision-making. The previous two chapters provided the necessary background in each of these areas individually. This chapter surveys the works that combine them, systems that use reinforcement learning to automate compiler optimization decisions.

The landscape of \gls{rl}-based compiler optimization can be organized into three tiers. First, \gls{rl} has been integrated directly into existing compiler frameworks such as Halide and Tiramisu, where the learning agent replaces or augments hand-crafted scheduling heuristics. Second, general-purpose \gls{rl} environments for compiler research have been developed, with CompilerGym and PolyGym providing standardized benchmarks and interfaces. Third, within the \gls{mlir} ecosystem, work initiated by Bendib et al.~\cite{bendib2024reinforcement} and substantially extended by Tirichine et al.~\cite{tirichine2025reinforcement} has established a growing line of \gls{rl}-based auto-schedulers for \gls{mlir} Linalg operations and full programs.

\section{\gls{rl} in Domain-Specific Compilers}

\subsection{Halide}

Halide is a domain-specific language and compiler for image processing and computational photography, introduced by Ragan-Kelley et al.~\cite{ragan2013halide}. A key innovation of Halide is the separation of the algorithm (what to compute) from the schedule (how to compute it). The algorithm is expressed as a functional description of the computation that operates on multi-dimensional arrays, while the schedule specifies loop transformations, tiling sizes, vectorization widths, parallelization granularity, fusion decisions, and storage layout, all through a small set of declarative scheduling directives.

This separation makes Halide an attractive target for automated scheduling, because the schedule is a structured, parameterized object that an optimization algorithm can search over independently of algorithmic correctness. The original Halide auto-scheduler used a beam search guided by a hand-crafted cost model that estimated execution time based on arithmetic intensity and memory access patterns. The search explored a tree of scheduling decisions, evaluating partial schedules with the cost model and pruning branches with low estimated performance.

Building on this architecture, Pecenin et al.~\cite{pecenin2019halide} replaced the beam search with a deep reinforcement learning approach. Their agent was trained to select scheduling directives sequentially for a given Halide pipeline. The state representation encoded loop nest features including loop bounds, buffer sizes, and the current partial schedule. The action space consisted of the same scheduling primitives available in the Halide scheduling language: tile, vectorize, parallelize, fuse, reorder, and compute-at directives with their associated parameters. The reward was the execution-time speedup relative to the unoptimized baseline.

The agent was trained using \gls{ppo}~\cite{schulman2017ppo} and achieved speedups competitive with the heuristic auto-scheduler, while requiring no manual tuning of the search strategy or hand-crafted cost model features. This result demonstrated that \gls{rl} can learn effective scheduling policies for domain-specific compilers.

\textbf{Limitations.} Halide's action space and cost model are tailored to image processing pipelines, which have a regular structure of stencil computations, upsampling, and downsampling. The scheduling primitives are specific to this domain and do not directly transfer to the broader set of linear algebra operations found in deep learning workloads. Furthermore, Halide is a standalone compiler rather than part of a general-purpose compiler infrastructure, limiting the applicability of its auto-scheduling techniques to other \glspl{ir} and optimization domains. The agent also does not incorporate hardware-specific features into its observation, making its decisions implicitly tied to the target platform used during training.

\subsection{Tiramisu}

Tiramisu is a polyhedral compiler designed for expressing and optimizing loop nests, introduced by Baghdadi et al.~\cite{baghdadi2018tiramisu}. Like Halide, Tiramisu separates the algorithm from the schedule and provides a scheduling language for specifying loop transformations. Tiramisu uses the polyhedral model to represent loop nests and their dependencies, enabling precise dependence analysis and the application of complex composite transformations. The polyhedral model represents loop iterations as integer points within a multi-dimensional polyhedron, and transformations correspond to linear mappings on this iteration space.

Within the same research group, work by Lamouri and Merad~\cite{lamouri2024tiramisu} explored the use of reinforcement learning to automate schedule selection in Tiramisu. Their agent learned to select sequences of scheduling commands (tiling, interchange, parallelization, vectorization, unrolling) for polyhedral loop nests. The state representation used features extracted from the polyhedral representation: iteration domain sizes, dependence vectors, access functions, and loop nest depth.

The \gls{mdp} formulation treated each scheduling command as an action that transitions the program from one partial schedule to another. The reward was the execution-time speedup achieved by the final schedule. The policy network was an \gls{lstm}-based architecture that processed the sequence of features extracted from the polyhedral representation.

An earlier effort by Hennouni and El Hassane~\cite{henouni2022utilisation} had explored \gls{rl} for Tiramisu using a feedforward policy network with a fully observed \gls{mdp}, where the entire state was visible at every step. Lamouri and Merad improved on this by adopting a partially observable \gls{mdp} (\gls{pomdp}) formulation with an \gls{lstm} encoder, which allowed the agent to accumulate knowledge about the program structure across multiple observation steps.

\textbf{Limitations.} The polyhedral model imposes restrictions on the programs that can be represented: loop bounds and array indices must be affine functions of enclosing loop iterators and symbolic constants. Many operations in modern machine learning workloads, particularly those with dynamic tensor shapes, indirect memory accesses, or non-affine control flow, fall outside these constraints. Operations such as convolutions with padding, depthwise separable convolutions, and attention mechanisms often violate the affine restriction. Additionally, Tiramisu is a research compiler with a narrower scope than \gls{mlir}, which is part of the \gls{llvm} ecosystem with broad industry adoption and an active development community.

\section{General-Purpose \gls{rl} Environments for Compilers}

\subsection{CompilerGym}

CompilerGym, introduced by Cummins et al.~\cite{cummins2021compilergym}, is a reinforcement learning environment for compiler optimization tasks. It provides a standardized OpenAI Gym interface, enabling researchers to apply \gls{rl} algorithms to production compiler passes without needing to implement the compiler interaction layer themselves.

CompilerGym exposes two main categories of benchmark tasks. The first is phase ordering, where the agent selects the sequence of \gls{llvm} optimization passes to apply to a program. The action space is a discrete set of \gls{llvm} passes such as \texttt{-mem2reg}, \texttt{-loop-unroll}, and \texttt{-inline}, and the agent decides the order and repetition of these passes. The second is flag tuning, where the agent selects compiler command-line flags to control optimization behavior globally.

The observation space consists of program features extracted by \gls{llvm}'s built-in analysis passes: instruction counts by opcode, basic block counts, loop depth histograms, and control-flow graph properties. The reward is typically the reduction in code size or the improvement in execution time relative to a baseline compilation using default optimization flags. CompilerGym ships with a large dataset of programs from diverse sources, including C benchmarks, C++ programs, and the \gls{llvm} test suite.

\textbf{Limitations.} CompilerGym's action space operates at the granularity of entire optimization passes applied to whole translation units, rather than per-operation or per-loop transformations. This makes it suitable for coarse-grained pass selection but limits its ability to express fine-grained schedules where different transformations are applied to different parts of the same program. The observation space, while rich in aggregate statistics, does not include the structured loop-level features (access matrices, loop bounds, iterator types) that are critical for loop nest optimization. Furthermore, the \gls{llvm}-level representation is lower-level than Linalg operations, making it harder to reason about high-level transformations such as tiling and fusion that are naturally expressed on structured linear algebra operations.

\subsection{PolyGym}

PolyGym, introduced by Brauckmann et al.~\cite{brauckmann2021polygym}, is a reinforcement learning environment specifically designed for polyhedral loop optimization. It builds on the Polyite polyhedral optimization framework and provides an \gls{rl} interface for selecting transformation sequences on polyhedral programs.

The PolyGym state representation encodes detailed polyhedral features of the loop nest: dependence polyhedra describing which iterations depend on which others, iteration space geometry including bounds and strides, and access functions mapping array subscripts to loop iterators. The action space includes polyhedral transformations such as tiling with parametric tile sizes, interchange, skewing, and parallelization, each with configurable parameters. The reward is the measured execution time of the transformed program on the target hardware.

PolyGym supports multiple benchmark suites, including PolyBench (a collection of kernels from linear algebra and scientific computing) and the Rodinia benchmark suite for heterogeneous computing. Its design enables head-to-head comparisons of different \gls{rl} algorithms on the same set of optimization problems.

\textbf{Limitations.} PolyGym shares the fundamental limitation of polyhedral methods: it can only represent programs with static, affine loop bounds and array accesses. This excludes many operations commonly found in machine learning compilers, such as convolutions with non-unit strides, depthwise separable convolutions, and attention mechanisms with dynamic shapes. Additionally, the polyhedral representation is opaque to the learning agent: the agent receives numerical features extracted from the polyhedra rather than the structured program graph or \gls{ast} that a compiler developer might reason about, making the learned policy difficult to interpret.

\section{\gls{rl} for \gls{mlir} Optimization}

This section provides an in-depth analysis of reinforcement learning applied to \gls{mlir}-based code optimization. The work in this area originates from a sustained research effort that began with the first \gls{rl} environment for single \gls{mlir} operations and progressed to systems capable of optimizing full computation graphs across multiple domains.

\subsection{Bendib et al. (2024): First \gls{rl} Environment for Single Operations}

The work by Bendib et al.~\cite{bendib2024reinforcement} established the first reinforcement learning environment for automatic code optimization in the \gls{mlir} compiler. Their system targets individual Linalg operations (matrix multiplications, convolutions, element-wise additions) and extracts structured numerical features from each operation to form the agent's observation.

\textbf{Environment design.} The environment defines a discrete, hierarchical action space. At the top level, the agent selects an action type from among Tiling (\texttt{T}), TiledParallelization (\texttt{TP}), Interchange (\texttt{I}), Vectorization (\texttt{V}), and NoTransformation (\texttt{NT}). Once the action type is chosen, the agent selects the parameters for that action through a type-specific parameter head. For tiling actions, the parameters are tile sizes for each loop in the nest, chosen from a set of powers of two (a tile size of zero indicates that loop should not be tiled). For interchange, the parameter is a permutation of the loop indices, encoded through one of three mechanisms: pre-enumerated candidate permutations, pointer-based sequential selection, or continuous Lehmer-code decoding. Vectorization is a terminal action: once applied, no further transformations can be applied to that operation.

The reward function is the log-speedup relative to the unoptimized baseline execution time. Penalties are applied for invalid transformations (e.g., tile sizes that do not divide loop bounds, a restriction specific to the \gls{mlir} \gls{rl} implementation rather than a general tiling constraint) and for execution failures during the compilation and \gls{jit} execution pipeline.

\textbf{State representation and policy.} The observation space encodes a fixed-size feature vector for each operation: operation type (matmul, conv, add, etc.), loop upper bounds, iterator types (parallel, reduction), load and store access matrices describing which loop iterators index which array dimensions, and arithmetic operation counts. These features are extracted directly from the \gls{mlir} Linalg operation using the \gls{mlir} Python bindings.

The policy network follows an actor-critic architecture. The operation features are processed through a multi-layer perceptron, producing an intermediate representation that is then passed through an \gls{lstm} to produce a fixed-size embedding vector encapsulating sequential dependencies across loop levels. A shared backbone network processes this embedding before branching into the action-type selection head, per-action parameter heads (one for each action type), and a value head that estimates the expected return from the current state. The \gls{ppo} algorithm~\cite{schulman2017ppo} is used for training, with Generalized Advantage Estimation~\cite{schulman2015gae} for computing advantage estimates.

\textbf{Contributions and limitations.} This work demonstrated that an \gls{rl} agent can learn effective optimization schedules for individual \gls{mlir} Linalg operations, establishing the viability of \gls{rl}-based optimization in the \gls{mlir} ecosystem. However, the system has several important limitations. The agent operates on single operations in isolation, without awareness of producer-consumer relationships or the effect of transformations on the broader computation graph. The observation includes no hardware information, preventing the agent from adapting its decisions to different target platforms. The reward is entirely sparse, received only at episode termination, which slows learning and provides no guidance on intermediate transformation choices. Finally, the \gls{lstm} encoder processes loop features sequentially from innermost to outermost, which may limit its ability to capture dependencies between non-adjacent loop levels in deeply nested structures.

\subsection{Tirichine et al. (2025): Extensions and Optimizations}

The work by Tirichine et al.~\cite{tirichine2025reinforcement} represents a substantial extension of the Bendib et al. baseline, advancing the system from single-operation optimization to full-program optimization with enhanced modeling capabilities. It advances the prior system with new methods for state encoding, action space design, and learning algorithms.

\textbf{\gls{ast}-based state representation and recursive \gls{lstm} encoder.} A central contribution of Tirichine et al. is the replacement of the flat per-operation feature vector with a structured \gls{ast}-based state representation. The computation graph is represented as a tree where leaf nodes correspond to individual Linalg operations and internal nodes encode the hierarchical structure of producer-consumer relationships. Each operation node carries the same features as in the Bendib baseline (operation type, loop bounds, iterator types, access matrices).

To encode this tree-structured observation into a fixed-size vector suitable for the policy network, the authors introduce a recursive \gls{lstm} encoder. Beginning at the leaf nodes, the features of each operation are embedded by a multi-layer perceptron. The embeddings of consumer operations are then combined with the embeddings of their producer operations through an \gls{lstm} that processes the \gls{ast} bottom-up, producing an aggregated embedding at the root that summarizes the entire computation graph. This recursive structure naturally captures the dependence chain: the encoding of a consumer reflects not only its own features but also those of all operations whose output it consumes.

\textbf{Fusion action and producer-consumer traversal.} To exploit the producer-consumer relationships captured by the \gls{ast}, Tirichine et al. introduce a new Fusion action, denoted \texttt{TiledFusion} (\texttt{TF}). This action merges a producer operation into its consumer's loop nest using \gls{mlir}'s tiling and fusion mechanisms, eliminating the intermediate buffer that would otherwise store the producer's output. When combined with tiling and parallelization, the action becomes \texttt{TiledParallelFusion} (\texttt{TPF}), which tiles both operations, fuses them, and parallelizes the resulting loop nest.

The environment traversal strategy is also redesigned to accommodate fusion decisions. Whereas the Bendib baseline processed operations in isolation, the Tirichine system traverses the computation graph from consumer to producer. This ordering ensures that when the agent visits a given operation, it has already observed and made optimization decisions for all operations that consume its output, providing critical context for fusion decisions. The agent can choose to fuse the current operation with a previously processed consumer, or optimize it independently.

\textbf{Action space innovations.} In addition to fusion, the work introduces two technical improvements to the interchange action. The original baseline offered three interchange encoding modes (enumerated candidates, pointer-based, and Lehmer-code continuous), each with distinct trade-offs between action space compactness and verbosity. Tirichine et al. introduce a pointer network architecture that directly selects a loop permutation from the set of all valid permutations using an attention mechanism over candidate loop positions, achieving more efficient exploration of the permutation space than the enumerated candidate approach while maintaining precision. A levels-pointer variant further improves on this by operating hierarchically, first selecting a coarse ordering and then refining it.

\textbf{Domain generalization: neural networks and LQCD.} A distinctive feature of the Tirichine system is its demonstration of generalization across two fundamentally different computational domains. The first domain consists of neural network operators common in deep learning: matrix multiplications, 2D convolutions, element-wise operations, and their compositions into blocks and full models (VGG, ResNet, MobileNet). The second domain is Lattice Quantum Chromodynamics (LQCD), a field of computational physics involving stencil-like computations on four-dimensional space-time grids.

To support LQCD workloads, the authors developed a dedicated compiler frontend that parses LQCD programs and lowers them to \gls{mlir} Linalg operations. This translation enables the same \gls{rl} auto-scheduler, trained on a mixture of neural network and LQCD programs, to optimize code from the physics domain without modification. The LQCD domain introduces computation patterns (stencils, reductions across multiple dimensions, complex index arithmetic) that differ markedly from the dense linear algebra patterns of neural networks, providing a stringent test of generalization.

\textbf{AlphaZero adaptation.} As an alternative to \gls{ppo}, the authors explore the AlphaZero algorithm~\cite{silver2017alphazero} for the code optimization problem. AlphaZero uses Monte Carlo Tree Search (MCTS) guided by a neural network to plan action sequences, potentially enabling deeper lookahead than the reactive policy of \gls{ppo}. In the adaptation, the MCTS simulates sequences of loop transformations from the current program state, with a learned value network evaluating partial sequences and a policy network proposing promising actions. The AlphaZero-based agent is compared against the \gls{ppo} baseline, revealing that while AlphaZero can discover higher-quality schedules in some settings, the computational cost of MCTS rollouts makes it less practical for large-scale training.

\textbf{Results and limitations.} The evaluation conducted by Tirichine et al. covers both single operations (various matrix multiplication and convolution sizes) and full neural network models (VGG, ResNet, MobileNet). On single operations, the \gls{rl} agent achieved speedups competitive with hand-tuned schedules. On full models, the agent demonstrated notable performance, including a 2.5$\times$ speedup on MobileNetV2 relative to standard PyTorch execution, though it underperformed on architectures with highly regular patterns (VGG, 1.09$\times$) and on ResNet (2.1$\times$ compared to PyTorch's 4.25$\times$).

Despite these advances, the authors identify several limitations that persist in their system. The observation remains hardware-agnostic: the agent must infer optimal tile sizes and parallelization strategies through trial and error rather than conditioning on known cache and core specifications. The reward function remains entirely sparse, with no intermediate signal guiding the agent toward promising states during an episode. The recursive \gls{lstm} encoder processes loop levels sequentially, which may struggle with deeply nested structures where the relationship between distant levels is not well captured. These limitations define the research space addressed by the present work.

\subsection{Other Efforts}

Beyond the USTHB line of research, the broader \gls{mlir} community has begun exploring learning-based optimization approaches, though published results remain limited. Hennouni and El Hassane~\cite{henouni2022utilisation} investigated the use of reinforcement learning for exploring the search space of \gls{mlir} transformations, focusing on loop nest optimization with a simplified action space. Their work confirmed the feasibility of \gls{rl}-guided optimization in \gls{mlir} but was limited to single-level loop nests and did not address the challenges of full computation graphs with producer-consumer dependencies. The \gls{mlir} project itself includes ongoing efforts to develop heuristic-based auto-scheduling capabilities through the transform dialect; these provide the infrastructure on which learning-based schedulers can operate, but the scheduling logic in these efforts remains heuristic-driven rather than learned.

\section{Comparative Synthesis}

This section synthesizes the systems surveyed above, identifying common design patterns, key differences, and the gaps that motivate the contributions of this thesis.

\subsection{Comparison Table}

\begin{table}[H]
    \centering
    \caption{Comparative synthesis of \gls{rl}-based compiler optimization systems.}
    \label{tab:related-work-synthesis}
    \begin{tabularx}{\textwidth}{l >{\raggedright\arraybackslash}p{2.5cm} >{\raggedright\arraybackslash}p{2.5cm} >{\raggedright\arraybackslash}p{2cm} >{\raggedright\arraybackslash}p{2cm} >{\raggedright\arraybackslash}p{1.5cm}}
        \toprule
        \textbf{System} & \textbf{Target Format} & \textbf{State Elements} & \textbf{Agent Architecture} & \textbf{Reward Type} & \textbf{HW-Aware} \\
        \midrule
        Halide + DRL~\cite{pecenin2019halide}
        & Halide
        & Loop nest features, buffer sizes
        & Feedforward / Beam Search
        & Sparse
        & No \\
        \addlinespace
        Tiramisu~\cite{lamouri2024tiramisu}
        & Polyhedral
        & Iteration domains, dependencies
        & \gls{lstm}
        & Sparse
        & No \\
        \addlinespace
        CompilerGym~\cite{cummins2021compilergym}
        & \gls{llvm} \gls{ir}
        & Aggregate program statistics
        & Feedforward
        & Sparse
        & No \\
        \addlinespace
        PolyGym~\cite{brauckmann2021polygym}
        & Polyhedral
        & Dependence polyhedra
        & Feedforward
        & Sparse
        & No \\
        \addlinespace
        \gls{mlir}-Base~\cite{bendib2024reinforcement}
        & \gls{mlir} Linalg
        & Flat op features
        & \gls{lstm}
        & Sparse
        & No \\
        \addlinespace
        \gls{mlir}-\gls{rl}~\cite{tirichine2025reinforcement}
        & \gls{mlir} Linalg
        & \gls{ast} block sequences
        & \gls{lstm}
        & Sparse
        & No \\
        \addlinespace
        \textbf{Proposed Agent}
        & \gls{mlir} Linalg
        & \textbf{\gls{ast} + HW profiles}
        & \textbf{Transformer}
        & \textbf{Shaped}
        & \textbf{Yes} \\
        \bottomrule
    \end{tabularx}
\end{table}

Table~\ref{tab:related-work-synthesis} summarizes the key design dimensions of the surveyed systems alongside our proposed system. Several patterns emerge from this comparison.

\subsection{Common Design Patterns}

All surveyed systems share a set of core design choices. Each models the optimization task as a Markov decision process where the agent sequentially selects transformations and receives execution-time-based feedback. The state representations are universally derived from features of the target \gls{ir}: loop nest dimensions, access patterns, operation types, and dependence information. The action spaces uniformly consist of compiler transformations with associated parameters, and all systems rely on execution time speedup as the primary or sole component of the reward function. \gls{ppo} is the most widely adopted learning algorithm among systems that train an agent, due to its stability, sample efficiency, and ability to handle mixed discrete-continuous action spaces.

\subsection{Limitations and Open Challenges}

The comparative analysis reveals several gaps in the existing literature that motivate the contributions of this thesis:

\begin{enumerate}
    \item \textbf{Sparse reward problem.} Every prior \gls{mlir}-based system relies exclusively on terminal rewards received at episode completion. In a typical episode where the agent applies 3 to 5 transformations before reaching a terminal action, only the final action receives a reward signal while all intermediate actions receive zero feedback. This creates a difficult credit assignment problem: the agent must infer which actions in the sequence contributed to the final outcome without any per-step guidance, slowing convergence and limiting the quality of the learned policy.

    \item \textbf{Hardware opacity.} None of the prior systems include hardware characteristics in the agent's observation. The agent must implicitly learn platform-appropriate strategies through execution-time feedback, which is both sample-inefficient (the agent must rediscover effective strategies for each platform) and non-transferable (a policy trained on one machine cannot adapt to a different machine without retraining). Encoding cache sizes, core counts, and \gls{simd} width directly into the observation would enable the policy to condition its decisions, particularly tile sizes and parallelization granularity, on the target platform.

    \item \textbf{Encoder expressiveness.} The recursive \gls{lstm} encoder used in prior \gls{mlir}-based work processes loop levels sequentially from innermost to outermost. While effective for capturing local dependencies between adjacent loop levels, this sequential processing may not adequately capture long-range interactions between non-adjacent levels, such as the relationship between the outermost tiling loop and the innermost vectorization loop. A transformer-based encoder with self-attention would allow every loop level to attend directly to every other level, regardless of their distance in the nest hierarchy.

    \item \textbf{Evaluation fragmentation.} Most prior work evaluates on a single benchmark suite or a small set of programs, often without standardized evaluation protocols or multi-configuration comparisons. This makes it difficult to isolate the contribution of individual design choices and to compare results across systems. A systematic evaluation framework with shared datasets, consistent training protocols, and controlled ablation studies is needed to advance the field.
\end{enumerate}

The contributions of this thesis directly address each of these challenges. The hardware-aware observation, shaped reward function, and transformer-based loop nest encoder are detailed in the proposed system design in Chapter~4. Our unified Proposed Agent combines all of these improvements to evaluate their cumulative effect on optimization performance, and isolates each contribution via a rigorous ablation study.

\section{Conclusion}

This chapter has surveyed the landscape of reinforcement learning applied to compiler optimization. We examined systems that integrate \gls{rl} into existing domain-specific compilers (Halide, Tiramisu), general-purpose \gls{rl} environments for compiler research (CompilerGym, PolyGym), and the specific line of work on \gls{rl}-based optimization in \gls{mlir} (Bendib et al., Tirichine et al., and related efforts).

The \gls{mlir}-based systems received in-depth treatment, reflecting their direct relevance to this thesis. Bendib et al. established the first \gls{rl} environment for single \gls{mlir} Linalg operations. Tirichine et al. extended this to full-program optimization with an \gls{ast}-based state representation, recursive \gls{lstm} encoder, fusion action, domain generalization across neural networks and LQCD, and algorithmic innovations including pointer networks for interchange and an AlphaZero adaptation.

The comparative synthesis revealed that while the surveyed systems share common design patterns (structured state representations, discrete transformation action spaces, execution-time-based rewards, and \gls{ppo} as the learning algorithm of choice), important gaps remain. The exclusive reliance on sparse terminal rewards, the absence of hardware awareness, the limitations of sequential \gls{lstm} encoders for hierarchical program structures, the restricted action diversity, and the fragmentation of evaluation protocols represent opportunities for meaningful improvement.

These gaps directly motivate the design decisions in our system. The next chapter presents the design and implementation of our \gls{rl}-based auto-scheduler, detailing the architectural enhancements, hardware-aware features, and shaped reward mechanisms that address these challenges.
